\documentclass[11pt,a4paper]{article}

\usepackage[a4paper,margin=0.78in]{geometry}
\usepackage[T1]{fontenc}
\usepackage[utf8]{inputenc}
\usepackage{lmodern}
\usepackage{microtype}
\usepackage{xcolor}
\usepackage{booktabs}
\usepackage{tabularx}
\usepackage{enumitem}
\usepackage{longtable}
\usepackage{hyperref}
\usepackage{fancyhdr}
\usepackage{titlesec}
\usepackage{array}

\definecolor{precedentgreen}{HTML}{2F6B45}
\definecolor{lightgreen}{HTML}{EEF6F0}
\definecolor{darkink}{HTML}{1E2933}
\definecolor{muted}{HTML}{5B6670}
\definecolor{warning}{HTML}{9A6700}

\hypersetup{
  colorlinks=true,
  linkcolor=precedentgreen,
  urlcolor=precedentgreen,
  citecolor=precedentgreen,
  pdftitle={Precedent - Updated Project Report},
  pdfauthor={We Code Coders}
}

\pagestyle{fancy}
\fancyhf{}
\lhead{\textcolor{precedentgreen}{\textbf{Precedent}}}
\rhead{\textcolor{muted}{Updated Project Report}}
\cfoot{\textcolor{muted}{\thepage}}
\setlength{\headheight}{14pt}

\titleformat{\section}
  {\Large\bfseries\color{precedentgreen}}
  {\thesection}{0.6em}{}
\titleformat{\subsection}
  {\large\bfseries\color{darkink}}
  {\thesubsection}{0.6em}{}

\newcommand{\status}[1]{\textcolor{precedentgreen}{\textbf{#1}}}
\newcommand{\limitation}[1]{\textcolor{warning}{\textbf{#1}}}
\newcommand{\product}{\textit{Precedent}}

\begin{document}

\begin{titlepage}
\thispagestyle{empty}
\vspace*{1.0cm}
\begin{center}
  {\Huge\bfseries\color{precedentgreen} Precedent\par}
  \vspace{0.25cm}
  {\Large\color{darkink} A Smarter Way to Manage Brand Rules and Approvals\par}
  \vspace{1.0cm}
  \rule{0.78\textwidth}{1pt}
  \vspace{0.65cm}

  {\Large\bfseries Updated Project Report\par}
  \vspace{0.35cm}
  {\large Inkloom AI Brand Governance Challenge\par}
  \vspace{1.2cm}

  \begin{tabular}{rl}
    \textbf{Team:} & We Code Coders\\
    \textbf{Team leader:} & Samarth Manoj Darak\\
    \textbf{Institution:} & Vishwakarma Institute of Pune\\
    \textbf{Product:} & Precedent\\
    \textbf{Report status:} & Final build / Round 2 completion\\
  \end{tabular}

  \vfill
  \begin{minipage}{0.78\textwidth}
  \centering
  \colorbox{lightgreen}{\parbox{0.92\linewidth}{
    \centering
    \textbf{Product status:} Core creator-to-manager governance workflow
    implemented, persisted in a database, and validated with automated tests.
  }}
  \end{minipage}
\end{center}
\end{titlepage}

\tableofcontents
\newpage

\section{Executive summary}

\product{} is an AI-assisted brand-governance and content-approval platform. It connects
creative content to structured brand rules, stored evidence, human approvals, scoped
exceptions, and a rule-change impact simulator.

The current build is deliberately human-controlled. Deterministic rules are the source of
truth. Optional AI helpers can suggest rules, semantic findings, explanations, and minimal
copy edits, but they cannot activate policy or approve content. When no LLM credentials
are configured, the application continues to work without fabricating AI output.

The verified workflow is:
\begin{enumerate}[leftmargin=*,itemsep=2pt]
  \item A creator signs in and submits copy or content with asset facts.
  \item The server evaluates the submission against active brand rules.
  \item Findings show severity, explanation, suggested correction, and source evidence.
  \item The creator can revise content, comment, or request a scoped exception.
  \item A manager approves, requests changes, or rejects a specific version.
  \item Managers can simulate proposed rule changes before activation.
  \item Decisions, exceptions, revisions, and changes remain available in the audit trail.
\end{enumerate}

\section{Completed and working features}

\begin{longtable}{p{0.32\textwidth}p{0.60\textwidth}}
\toprule
\textbf{Feature} & \textbf{Completed behaviour}\\
\midrule
\endhead
Authentication and RBAC & JWT authentication, Argon2 password hashing, role-aware routes, and server-enforced creator/manager permissions.\\
Guideline library & Manager upload of PDF, DOCX, and text-based guidelines; extraction of text and stored passages with page information.\\
Structured rules & Draft rules, human confirmation, active rules, rule metadata, effective dates, scopes, priorities, and source references.\\
Content review & Text and asset submissions evaluated by the deterministic server-side rule engine.\\
Evidence-linked findings & Findings carry rule codes, explanations, source documents, source passages, and evidence verification state.\\
Versioned approvals & Content versions, comments, manager decisions, and version-specific review history.\\
Scoped exceptions & Exception reason, scope, effective date, expiry date, conditions, manager decision, history, and revocation.\\
Expiry enforcement & Expired exceptions cannot be approved by the backend.\\
Conflict detection & Deterministic detection of conflicts between active rules with explainable results.\\
Rule Impact Sandbox & Proposed rule changes are replayed against stored submissions and reported as compliant, needs review, or non-compliant.\\
Audit and health & Append-only audit history, filtering, guideline health metrics, expiring exceptions, and CSV/JSON export.\\
Asset library & Brand/product asset upload, approval, attachment, and stored asset facts used in visual rule checks.\\
Responsive UI & Shared React/Tailwind design system with loading, empty, error, success, and permission states across creator and manager screens.\\
\bottomrule
\end{longtable}

\section{Technology and architecture}

\subsection{Application stack}

\begin{itemize}[leftmargin=*,itemsep=2pt]
  \item \textbf{Frontend:} React 18, TypeScript, Vite, Tailwind CSS, React Router, and reusable UI components.
  \item \textbf{Backend:} FastAPI, Pydantic, SQLAlchemy, and Python 3.12.
  \item \textbf{Database:} SQLite for the reproducible demo; PostgreSQL-ready through \texttt{DATABASE\_URL} and Alembic.
  \item \textbf{Authentication:} Argon2 password hashing and signed JWT tokens.
  \item \textbf{Document extraction:} PyMuPDF for PDF files and \texttt{python-docx} for DOCX files.
  \item \textbf{Optional AI:} OpenAI-compatible chat-completions endpoint with Pydantic-validated structured output.
\end{itemize}

\subsection{How the pieces fit together}

The React application communicates with the FastAPI application through one typed API
service layer. The frontend does not maintain a separate mock data store. Authenticated
reads and writes go to the backend, and mutations are persisted before the UI refreshes.

The backend separates routers, services, models, and the rule engine. Routers enforce
authentication and manager-only operations. Services handle extraction, submissions,
exceptions, simulations, assets, and audit records. SQLAlchemy models preserve the
relationship between submissions, versions, findings, source passages, decisions, and
exceptions.

\subsection{The technically difficult part}

The central technical challenge was converting natural-language brand guidance into
reviewable, testable rules without allowing an AI model to invent policy. The implementation
stores each rule with a code, subject, condition, action, scope, priority, source passage,
and effective dates. The deterministic engine evaluates explicit conditions and returns
\texttt{unknown} when the required asset fact is unavailable. This prevents missing
information from becoming a false violation.

Every finding is linked back to a stored guideline passage. Optional AI output is validated,
restricted to known rule codes where applicable, and treated as assistance rather than
authority. Exception decisions are also recorded as first-class, auditable objects instead
of being left in an informal comment or chat.

\section{Validation and current results}

\begin{table}[h]
\centering
\begin{tabularx}{\textwidth}{>{\bfseries}p{0.36\textwidth}X}
\toprule
\textbf{Check} & \textbf{Result}\\
\midrule
Backend automated tests & \status{41 tests passed} using the Python 3.12 environment.\\
Frontend typecheck and production build & \status{Passed} with \texttt{npm run build}.\\
Authentication and RBAC & \status{Passed} for creator and manager flows.\\
Exception expiry security & \status{Passed}; expired approval returns a conflict response.\\
Evidence verification & \status{Passed}; findings resolve to stored source passages.\\
Rule simulation & \status{Passed}; simulations evaluate stored submissions before activation.\\
Browser workflow & \status{Passed in the project QA report}; creator submission to manager queue and audit flow completed with no console errors.\\
\bottomrule
\end{tabularx}
\end{table}

\section{What changed from the original proposal}

The core direction did not change: the product still focuses on evidence-backed review,
human approvals, scoped exceptions, auditability, and pre-activation rule simulation.

SQLite became the default demonstration database so judges can run the system without
provisioning a separate database. PostgreSQL remains supported through configuration.
The planned pgvector/embedding path was not implemented because the deterministic,
evidence-linked workflow was more important to complete and validate in the available time.

The build also added an asset library, claims management, exception revocation, audit
export, workspace isolation, and server-side RBAC. OCR, advanced computer vision,
multilingual support, publishing integrations, and post-publication monitoring were kept
out of the current release rather than presenting partially working features as complete.

\section{Known limitations}

\begin{itemize}[leftmargin=*,itemsep=4pt]
  \item \limitation{Live LLM assistance is optional and not configured in the demo environment.}
        The deterministic engine remains fully usable without an API key.
  \item \limitation{Embedding-based retrieval with pgvector is not implemented.}
  \item \limitation{Full OCR and automatic visual understanding of uploaded posters are not implemented.}
        Visual checks currently use recorded asset metadata.
  \item \limitation{Notifications and post-publication monitoring are not implemented.}
  \item \limitation{The pinned backend dependency set requires Python 3.12.}
  \item The optional external image-source provider is not configured.
\end{itemize}

\section{Round 3: three major improvements}

\subsection{1. Full visual and OCR compliance analysis}

Add image processing and OCR so creators can upload posters and social creatives directly.
The system will inspect visible text, logo placement, colours, typefaces, clear space,
aspect ratio, and unreadable regions instead of relying primarily on manually entered
asset metadata.

\subsection{2. Evidence-grounded semantic AI review}

Add embedding-based retrieval and enable the configured LLM workflow for long guideline
documents and nuanced copy. Retrieved passages will be shown beside findings, and every
AI citation will be checked against stored evidence. Deterministic violations, AI
suggestions, and human-review cases will remain visibly distinct.

\subsection{3. Post-publication Sentinel and change-impact notifications}

Add a monitoring workflow for approved content and policy changes. When an approved claim,
exception, or rule changes, the system will identify affected content, notify responsible
users, and create a reassessment task. This extends governance beyond the initial approval.

\section{Demonstration flow}

\begin{enumerate}[leftmargin=*,itemsep=3pt]
  \item Sign in as the content creator.
  \item Submit sample copy such as ``100\% safe, boosts immunity, 20\% off''.
  \item Open the findings and inspect rule codes, explanations, and source evidence.
  \item Request an exception for a finding.
  \item Sign in as the brand manager and open the review queue.
  \item Approve, reject, or request changes for the exact submitted version.
  \item Open Exceptions to show scoped decisions and expiry enforcement.
  \item Open Rule Impact Sandbox to simulate a proposed policy change.
  \item Open Audit History to verify the recorded workflow.
\end{enumerate}

\section{Reproducibility}

\begin{verbatim}
Frontend:
  npm install
  npm run dev

Backend (Python 3.12):
  cd backend
  python -m venv .venv
  .venv\Scripts\activate
  pip install -r requirements.txt
  python seed.py
  uvicorn app.main:app --reload --port 8000

Validation:
  cd backend
  pytest -q
  cd ..
  npm run build
\end{verbatim}

\section{Declaration}

This product report describes the implementation in the repository at the time of
submission. The demonstration uses fictional brand data. The team understands that
production deployment for regulated or sensitive organisations would require additional
security review, policy controls, privacy review, and specialist legal oversight.

\vfill
\begin{center}
\textcolor{muted}{\small Prepared for the Inkloom AI Brand Governance Challenge}
\end{center}

\end{document}
